% GENERATED FILE. Edit canonical lessons, then run npm run generate:books.
% canonical-source-hash: fnv1a64:4ee7ab3ce32f6e2a
% canonical-lessons: LA-C333-festus, LA-C333-sollemnis, LA-C333-pompa, LA-C333-epulae, LA-C333-epulor

\chapter{Diēs fēstus}
\label{ch:la-dies-festus}
\hlchaptermodality{333}

\begin{chapteropening}
\textbf{By the end of this chapter:} \emph{I can describe a festival day: festive, yearly and solemn, a procession, a feast and to feast.}

Complete the last lesson of chapter 333: I can describe a festival day: festive, yearly and solemn, a procession, a feast and to feast.
\end{chapteropening}

\section[fēstus, fēsta, fēstum]{fēstus, fēsta, fēstum — festive, holiday}

\label{lesson:LA-C333-festus}



\begin{quote}
Before the new one: say the Latin for to seal, then the Latin for to leave behind.
\end{quote}

\subsection*{You'll want to know: fēstus, fēsta, fēstum}
\textbf{fēstus, fēsta, fēstum} — \textquotedblleft{}festive, holiday\textquotedblright{}.

\begin{quote}
Diē fēstō nēmō in forō labōrābat.
\end{quote}

\begin{quote}
\emph{On the holiday nobody was working in the forum.}
\end{quote}

\emph{Diēs fēstus}: a holiday, the opposite of a working day. English \emph{feast} and \emph{festival} come from it.

\begin{cousinweb}
\textbf{fēstus} belongs, says Lewis \& Short, with \emph{fēriae}, \textquotedblleft{}holidays\textquotedblright{} (taught). English \emph{feast} comes from it through French, and \emph{festival} from its family.
\end{cousinweb}

\begin{grammarlens}[title={What you've built}]
One more word for telling, explaining and describing at B1.
\end{grammarlens}

\subsection*{Your turn}
\begin{itemize}
\raggedright
  \item \emph{Say it:} \emph{fēstus, fēsta, fēstum}
  \item \emph{Say it:} \emph{fēstus, fēsta, fēstum}, once more
  \item \emph{Say it:} a sentence of your own with \emph{fēstus, fēsta, fēstum}
  \item \emph{Recall:} say the Latin for to seal, then the Latin for to leave behind, then say \emph{fēstus, fēsta, fēstum} again
\end{itemize}

\begin{tcolorbox}[breakable,colback=teal!4,colframe=teal!35!black,title={Before you move on}]
What does fēstus, fēsta, fēstum mean? (\textquotedblleft{}Festive, holiday\textquotedblright{}.) Say it once more.
\end{tcolorbox}
\section[sollemnis, sollemne]{sollemnis, sollemne — yearly, solemn; customary}

\label{lesson:LA-C333-sollemnis}



\begin{quote}
Before the new one: say the Latin for to leave behind, then the Latin for festive.
\end{quote}

\subsection*{You'll want to know: sollemnis, sollemne}
\textbf{sollemnis, sollemne} — \textquotedblleft{}yearly, solemn; customary\textquotedblright{}.

\begin{quote}
Sacrificium sollemne quotannīs fīēbat.
\end{quote}

\begin{quote}
\emph{A solemn sacrifice was held every year.}
\end{quote}

Two endings. A rite that comes round every year, and so solemn and customary. English \emph{solemn}.

\begin{cousinweb}
\textbf{sollemnis} is explained by Lewis \& Short as \emph{sollus}, \textquotedblleft{}whole\textquotedblright{}, and \emph{annus}, \textquotedblleft{}year\textquotedblright{}: what happens every year; EtymDB derives it from \emph{sollus}. English \emph{solemn} comes from it through French.
\end{cousinweb}

\begin{grammarlens}[title={What you've built}]
One more word for telling, explaining and describing at B1.
\end{grammarlens}

\subsection*{Your turn}
\begin{itemize}
\raggedright
  \item \emph{Say it:} \emph{sollemnis, sollemne}
  \item \emph{Say it:} \emph{sollemnis, sollemne}, once more
  \item \emph{Say it:} a sentence of your own with \emph{sollemnis, sollemne}
  \item \emph{Recall:} say the Latin for to leave behind, then the Latin for festive, then say \emph{sollemnis, sollemne} again
\end{itemize}

\begin{tcolorbox}[breakable,colback=teal!4,colframe=teal!35!black,title={Before you move on}]
What does sollemnis, sollemne mean? (\textquotedblleft{}Yearly, solemn; customary\textquotedblright{}.) Say it once more.
\end{tcolorbox}
\section[pompa, pompae]{pompa, pompae — a procession}

\label{lesson:LA-C333-pompa}



\begin{quote}
Before the new one: say the Latin for festive, then the Latin for yearly.
\end{quote}

\subsection*{You'll want to know: pompa, pompae}
\textbf{pompa, pompae} — \textquotedblleft{}a procession\textquotedblright{}.

\begin{quote}
Longa pompa per viās urbis prōcessit.
\end{quote}

\begin{quote}
\emph{A long procession moved through the streets of the city.}
\end{quote}

Feminine. Games, triumphs, weddings and funerals all had their \emph{pompa}. English \emph{pomp}.

\begin{cousinweb}
\textbf{pompa} was borrowed from Greek \emph{pompē}, \textquotedblleft{}a sending, a procession\textquotedblright{}, from \emph{pempō}, \textquotedblleft{}to send\textquotedblright{} (EtymDB). English \emph{pomp} comes from it through French.
\end{cousinweb}

\begin{grammarlens}[title={What you've built}]
One more word for telling, explaining and describing at B1.
\end{grammarlens}

\subsection*{Your turn}
\begin{itemize}
\raggedright
  \item \emph{Say it:} \emph{pompa, pompae}
  \item \emph{Say it:} \emph{pompa, pompae}, once more
  \item \emph{Say it:} a sentence of your own with \emph{pompa, pompae}
  \item \emph{Recall:} say the Latin for festive, then the Latin for yearly, then say \emph{pompa, pompae} again
\end{itemize}

\begin{tcolorbox}[breakable,colback=teal!4,colframe=teal!35!black,title={Before you move on}]
What does pompa, pompae mean? (\textquotedblleft{}A procession\textquotedblright{}.) Say it once more.
\end{tcolorbox}
\section[epulae, epulārum]{epulae, epulārum — a feast, a banquet}

\label{lesson:LA-C333-epulae}



\begin{quote}
Before the new one: say the Latin for yearly, then the Latin for a procession.
\end{quote}

\subsection*{You'll want to know: epulae, epulārum}
\textbf{epulae, epulārum} — \textquotedblleft{}a feast, a banquet\textquotedblright{}.

\begin{quote}
Post sacrificium epulae magnae parātae sunt.
\end{quote}

\begin{quote}
\emph{After the sacrifice a great feast was prepared.}
\end{quote}

Feminine, plural; the singular is the neuter \emph{epulum}, a public banquet.

\begin{cousinweb}
\textbf{epulae} is the plural that goes with \emph{epulum}, \textquotedblleft{}a banquet\textquotedblright{}; Lewis \& Short call its origin doubtful and suggest, with hesitation, a link with \emph{edō}, \textquotedblleft{}to eat\textquotedblright{}.
\end{cousinweb}

\begin{grammarlens}[title={What you've built}]
One more word for telling, explaining and describing at B1.
\end{grammarlens}

\subsection*{Your turn}
\begin{itemize}
\raggedright
  \item \emph{Say it:} \emph{epulae, epulārum}
  \item \emph{Say it:} \emph{epulae, epulārum}, once more
  \item \emph{Say it:} a sentence of your own with \emph{epulae, epulārum}
  \item \emph{Recall:} say the Latin for yearly, then the Latin for a procession, then say \emph{epulae, epulārum} again
\end{itemize}

\begin{tcolorbox}[breakable,colback=teal!4,colframe=teal!35!black,title={Before you move on}]
What does epulae, epulārum mean? (\textquotedblleft{}A feast, a banquet\textquotedblright{}.) Say it once more.
\end{tcolorbox}
\section[epulor, epulārī]{epulor, epulārī — to feast, to banquet}

\label{lesson:LA-C333-epulor}



\begin{quote}
Before the new one: say the Latin for a procession, then the Latin for a feast.
\end{quote}

\subsection*{You'll want to know: epulor, epulārī}
\textbf{epulor, epulārī} — \textquotedblleft{}to feast, to banquet\textquotedblright{}.

\begin{quote}
Tōtam noctem amīcī epulābantur.
\end{quote}

\begin{quote}
\emph{The friends feasted all night long.}
\end{quote}

\emph{epulor, epulārī, epulātus sum}, a deponent.

\begin{cousinweb}
\textbf{epulor} is built on \emph{epulum}, \textquotedblleft{}a banquet\textquotedblright{} (Lewis \& Short): to hold a feast.
\end{cousinweb}

\begin{grammarlens}[title={What you've built}]
One more word for telling, explaining and describing at B1.
\end{grammarlens}

\subsection*{Your turn}
\begin{itemize}
\raggedright
  \item \emph{Say it:} \emph{epulor, epulārī}
  \item \emph{Say it:} \emph{epulor, epulārī}, once more
  \item \emph{Say it:} a sentence of your own with \emph{epulor, epulārī}
  \item \emph{Recall:} say the Latin for a procession, then the Latin for a feast, then say \emph{epulor, epulārī} again
\end{itemize}

\begin{tcolorbox}[breakable,colback=teal!4,colframe=teal!35!black,title={Before you move on}]
What does epulor, epulārī mean? (\textquotedblleft{}To feast, to banquet\textquotedblright{}.) Say it once more.
\end{tcolorbox}
